\documentclass[11pt]{article}

\usepackage[margin=1in]{geometry}
\usepackage{enumitem}
\usepackage{amsmath, amssymb}

\setlist[enumerate]{itemsep=0.6em, topsep=0.6em}

\begin{document}

\begin{center}
  {\Large \textbf{PHIL 340 Examination}}\\[0.25em]
  {\large \textbf{Comparative Philosophy: East and West -- Answer Key}}\\[0.5em]
\end{center}

\begin{enumerate}[label=\textbf{\arabic*.}]

  \item \textbf{(B)} The action can be willed as a universal law without contradiction.\\
  Kant's Categorical Imperative requires that moral maxims be universalizable. An action is 
  permissible only if its underlying principle could be willed as a law for all rational beings 
  without logical contradiction or undermining itself.

  \item \textbf{(B)} \textit{Wu wei} (effortless action or non-interference).\\
  \textit{Wu wei} is central to Daoism, emphasizing spontaneous, natural action aligned with the 
  Dao (the Way). It suggests acting without forcing, allowing things to unfold naturally rather 
  than imposing artificial structures.

  \item \textbf{(B)} A disposition to act according to the mean between extremes.\\
  For Aristotle, virtue (\textit{arete}) is a state of character concerned with choice, lying in 
  a mean between two vices (excess and deficiency). For example, courage is the mean between 
  cowardice and recklessness.

  \item \textbf{(B)} \textit{Anatta} (non-self).\\
  \textit{Anatta} is the Buddhist doctrine that there is no permanent, unchanging self or soul. 
  All phenomena, including persons, are composite, impermanent, and interdependent. This is one 
  of the Three Marks of Existence in Buddhism.

  \item \emph{Sample answer:}\\
  \textbf{Similarities}:
  \begin{itemize}
    \item Both emphasize duty and moral obligation over consequences
    \item Both stress the importance of cultivating moral character
    \item Both reject pure consequentialism
    \item Both value universality (Kant's maxims, Confucius's reciprocity)
  \end{itemize}

  \textbf{Differences}:
  \begin{itemize}
    \item \textbf{Foundation}: Kant grounds morality in autonomous reason and categorical imperatives; 
    Confucius grounds it in social relationships and tradition (\textit{li}, ritual propriety)
    \item \textbf{Focus}: Kant emphasizes abstract, universal principles independent of context; 
    Confucianism emphasizes particular relationships (filial piety, loyalty) and social harmony
    \item \textbf{Method}: Kant uses rational formalism; Confucianism uses exemplars and cultivation 
    of virtues like \textit{ren} (benevolence)
    \item \textbf{Application}: Kant seeks absolute moral laws; Confucianism allows flexibility 
    based on relationships and circumstances
  \end{itemize}

  \item \emph{Sample answer:}\\
  \textbf{Atman in Hindu philosophy}:
  \begin{itemize}
    \item \textit{Atman} is the eternal, unchanging self or soul that persists through reincarnation
    \item It is distinct from the physical body and empirical ego
    \item The ultimate goal in Hinduism is realizing that \textit{Atman} is identical with \textit{Brahman}, 
    the universal, absolute reality: ``Tat Tvam Asi'' (That Thou Art)
    \item This realization leads to \textit{moksha} (liberation from the cycle of rebirth)
  \end{itemize}

  \textbf{Contrast with Buddhist Anatta}:
  \begin{itemize}
    \item Buddhism explicitly rejects the existence of \textit{Atman}
    \item According to \textit{Anatta}, there is no permanent self; what we call ``self'' is a 
    collection of five aggregates (\textit{skandhas}): form, sensation, perception, mental formations, 
    consciousness
    \item These aggregates are impermanent and constantly changing
    \item Liberation (\textit{nirvana}) comes from realizing this non-self and letting go of attachment 
    to the illusion of a permanent identity
  \end{itemize}

  \item \emph{Sample answer:}\\
  \textbf{Plato's Theory of Forms}: The physical world is a realm of imperfect copies; true reality 
  consists of eternal, immutable Forms (or Ideas) existing in a transcendent realm.

  \textbf{Epistemological implications}:
  \begin{itemize}
    \item True knowledge (\textit{episteme}) is of Forms, not sensible particulars
    \item Sensory experience provides only opinion (\textit{doxa})
    \item Knowledge is \textit{anamnesis} (recollection): the soul knew the Forms before birth and 
    recalls them through philosophical inquiry
  \end{itemize}

  \textbf{Metaphysical implications}:
  \begin{itemize}
    \item Addresses the problem of universals: what all triangles share is participation in the Form 
    of Triangle
    \item Forms are causally explanatory: things are what they are by participating in Forms
    \item The Form of the Good is the highest Form, source of truth and reality (Allegory of the Cave)
  \end{itemize}

  \textbf{Problem of universals}: How can many particular things share the same property? Plato answers: 
  by participating in a single universal Form that exists independently.

  \item \emph{Sample answer:}\\
  \textbf{Radical freedom}: Sartre argues that humans are ``condemned to be free.'' We have no fixed 
  nature or essence; existence precedes essence. We constantly create ourselves through choices.

  \textbf{Authenticity vs. Bad faith}:
  \begin{itemize}
    \item \textbf{Authenticity}: Acknowledging radical freedom, taking responsibility for choices, 
    living consciously without self-deception
    \item \textbf{Bad faith}: Self-deception where one denies freedom by pretending to be determined 
    by external factors (social roles, past, others' expectations). Example: a waiter who identifies 
    completely with his role, denying his freedom to be otherwise
  \end{itemize}

  \textbf{Moral responsibility}:
  \begin{itemize}
    \item With radical freedom comes absolute responsibility: we cannot blame circumstances, upbringing, 
    or nature for our choices
    \item We create values through our choices; there are no objective moral facts
    \item This can lead to existential anguish (\textit{angoisse}) as we confront the weight of 
    our freedom
  \end{itemize}

  Sartre's existentialism challenges us to live authentically by embracing our freedom and accepting 
  full responsibility for our lives.

\end{enumerate}

\end{document}
